% 例 1：长方体 4×2×3 的三视图（带尺寸）
\documentclass[tikz,border=4pt]{standalone}
\usepackage{ctex}
\usepackage{amsmath}
\usetikzlibrary{calc}
\begin{document}
\begin{tikzpicture}[thick, scale=0.7]
  % 主视：4×3
  \draw (0,0) rectangle (4,3);
  \node[font=\footnotesize] at (2,1.5) {主视图};
  \node[font=\footnotesize, below] at (2,0) {$4$};
  \node[font=\footnotesize, left] at (0,1.5) {$3$};
  % 左视：2×3
  \draw (5,0) rectangle (7,3);
  \node[font=\footnotesize] at (6,1.5) {左视图};
  \node[font=\footnotesize, below] at (6,0) {$2$};
  \node[font=\footnotesize, right] at (7,1.5) {$3$};
  % 俯视：4×2
  \draw (0,-3) rectangle (4,-1);
  \node[font=\footnotesize] at (2,-2) {俯视图};
  \node[font=\footnotesize, below] at (2,-3) {$4$};
  \node[font=\footnotesize, left] at (0,-2) {$2$};
  % 对齐虚线
  \draw[dashed, gray, very thin] (0,0) -- (0,-1);
  \draw[dashed, gray, very thin] (4,0) -- (4,-1);
  \draw[dashed, gray, very thin] (4,0) -- (5,0);
  \draw[dashed, gray, very thin] (4,3) -- (5,3);
\end{tikzpicture}
\end{document}
